\section{АНАЛИТИЧЕСКАЯ ЧАСТЬ}

\subsection{Нагрузочное тестирование СУБД: понятия, цели, методологии}

Производительность баз данных является одним из определяющих факторов качества современных информационных систем. С ростом объёмов обрабатываемых данных и увеличением числа одновременных пользователей задача объективной оценки производительности СУБД приобретает критическое значение. Нагрузочное тестирование представляет собой основной метод её решения, позволяя определить поведение системы под контролируемой нагрузкой и выявить потенциальные узкие места до их проявления в промышленной эксплуатации~\cite{gregg_2020}.

Нагрузочным тестированием называется вид тестирования производительности, направленный на определение поведения системы при ожидаемой и повышенной нагрузке, создаваемой путём эмуляции работы множества одновременных пользователей. Важно разграничивать нагрузочное тестирование и стресс-тестирование: первое оценивает работу системы в рамках проектных параметров нагрузки, второе определяет границы работоспособности при экстремальных значениях, превышающих нормальные эксплуатационные характеристики. Отдельным видом является тестирование стабильности , при котором система подвергается длительной нагрузке для выявления деградации производительности, утечек памяти и накопления фрагментации данных. Совместное применение указанных видов тестирования позволяет получить полную картину поведения СУБД в различных нагрузочных сценариях.

Ключевыми метриками, измеряемыми в процессе нагрузочного тестирования СУБД, являются задержка и пропускная способность~\cite{kleppmann_2017}. Задержка характеризует время от отправки запроса к базе данных до получения ответа и измеряется в миллисекундах. Для полного описания распределения задержек применяются процентили: медиана (p50) отражает типичное время отклика, p95 показывает задержку, которую не превышают 95\,\% запросов, а p99 характеризует наихудшие случаи. Использование исключительно среднего значения задержки является недостаточным, поскольку оно маскирует выбросы, способные существенно влиять на пользовательский опыт. Пропускная способность измеряется количеством успешно обработанных единиц нагрузки в секунду: запросов для SQL-сценариев и транзакций для транзакционных сценариев. Зависимость между задержкой и пропускной способностью нелинейна: при достижении определённого порога нагрузки пропускная способность перестаёт расти, а задержка начинает экспоненциально увеличиваться. Точка, в которой пропускная способность достигает максимума, называется точкой насыщения и является важным параметром при определении предельной нагрузки на систему.

Помимо основных метрик, в процессе тестирования отслеживаются системные показатели: загрузка центрального процессора, использование оперативной памяти, интенсивность дискового ввода-вывода. Из внутренних показателей СУБД особый интерес представляют коэффициент попаданий в кэш и количество взаимных блокировок. Коэффициент попаданий в кэш отражает долю запросов, обработанных без обращения к диску, и служит индикатором эффективности использования буферного пула. Взаимные блокировки возникают при конкурентном доступе к данным и приводят к принудительному завершению части транзакций, снижая общую пропускную способность. Рост частоты блокировок при увеличении числа параллельных транзакций сигнализирует о необходимости оптимизации структуры запросов или уровня изоляции транзакций.

Стандартизацией методологий тестирования производительности СУБД занимается организация TPC~--- всемирный консорциум, устанавливающий стандарты для оценки быстродействия и надёжности систем обработки транзакций и баз данных~\cite{tpc_specs}. Наиболее значимыми спецификациями TPC для реляционных СУБД являются TPC-C, TPC-H и TPC-E.

Спецификация TPC-C моделирует среду обработки заказов, в которой множество пользователей одновременно выполняют транзакции нескольких типов: оформление заказов, регистрация платежей, проверка статусов и мониторинг складских запасов. Бенчмарк охватывает пять типов параллельных транзакций различной сложности, обеспечивая комплексную проверку компонентов СУБД. Основной метрикой TPC-C является количество новых заказов в минуту~--- tpmC. Данная спецификация широко применяется для оценки OLTP-систем и реализована в ряде инструментов тестирования, в частности в HammerDB под наименованием TPROC-C. Спецификация TPC-H ориентирована на системы поддержки принятия решений и представляет собой набор бизнес-ориентированных ad-hoc запросов и модификаций данных; её основная метрика~--- составной показатель производительности запросов в час QphH@Size; в инструментах тестирования данная нагрузка фигурирует под наименованием TPROC-H. Спецификация TPC-E представляет более современную OLTP-модель, воспроизводящую деятельность брокерской фирмы с асинхронными операциями и сложным профилем чтения и записи; метрика~--- количество транзакций в секунду. По сравнению с TPC-C, TPC-E обеспечивает более реалистичную нагрузку для современных онлайн-систем.

Параллельно с автоматизацией тестирования сформировалась концепция наблюдаемости, объединяющая сбор, хранение и анализ телеметрических данных в режиме реального времени. Ключевым компонентом такого стека в области мониторинга метрик является система Prometheus~--- открытое программное обеспечение для сбора и хранения временных рядов с многомерной моделью данных~\cite{prometheus_overview}. Данные в Prometheus идентифицируются по имени метрики и набору меток в формате ключ--значение, что позволяет выполнять гибкую фильтрацию и агрегацию с помощью языка запросов PromQL. Для унификации инструментирования приложений разработан стандарт OpenTelemetry~--- независимый от поставщика фреймворк, охватывающий трассировки, метрики и журналы и обеспечивающий передачу телеметрии через единый компонент OpenTelemetry Collector в различные бэкенды~\cite{opentelemetry}. Совместное применение Prometheus и OpenTelemetry формирует полноценный стек наблюдаемости, позволяющий коррелировать результаты нагрузочных тестов с системными показателями. Необходимо отметить, что организация TPC расширяет область стандартизации за рамки реляционных СУБД: в частности, спецификация TPCx-AI адресована платформам машинного обучения и аналитики данных, однако данное направление выходит за рамки предметной области настоящей работы.

Типовая методология нагрузочного тестирования СУБД включает несколько последовательных этапов. На этапе подготовки определяется цель тестирования, выбирается или формируется тестовая база данных, создаются сценарии нагрузки и настраивается тестовый стенд. На этапе прогрева  система выводится на стабильный режим работы, исключающий влияние начальной инициализации кэшей и буферов. Основной этап заключается в выполнении нагрузки с заданным числом виртуальных пользователей и фиксацией метрик через равные интервалы времени. Завершающий этап анализа предполагает вычисление агрегированных показателей, построение зависимостей метрик от нагрузки и формулировку выводов.

Таким образом, нагрузочное тестирование СУБД представляет собой комплексный процесс, охватывающий определение целей, выбор методологии, генерацию контролируемой нагрузки, сбор и анализ метрик производительности. Стандартизированные спецификации TPC обеспечивают сопоставимость результатов между различными системами. Рассмотренные методологии и метрики составляют теоретическую основу для анализа существующих инструментальных решений, представленного в следующем разделе.

\subsection{Обзор существующих инструментов}

Прежде чем анализировать конкретные инструменты, необходимо определить критерии их отбора и логику классификации. В рамках настоящего обзора рассматриваются решения, непосредственно применимые к нагрузочному тестированию реляционных СУБД: инструменты должны поддерживать хотя бы одну из целевых СУБД, генерировать контролируемую нагрузку и собирать метрики производительности. По характеру своего происхождения и целевой аудитории все рассматриваемые решения разбиваются на три группы: классические инструменты для реляционных СУБД с открытым исходным кодом, современный DevOps-ориентированный инструментарий и коммерческие корпоративные решения. Такая классификация позволяет выявить не просто отдельные достоинства и недостатки продуктов, но закономерности, характерные для каждого класса в целом.

\subsubsection{Классические инструменты для реляционных СУБД}

Sysbench~--- кроссплатформенный инструмент нагрузочного тестирования и бенчмаркинга с открытым исходным кодом, изначально созданный для MySQL и впоследствии расширенный поддержкой PostgreSQL. Инструмент предоставляет набор встроенных OLTP-сценариев, охватывающих операции чтения, записи, обновления и удаления данных, и допускает создание пользовательских сценариев на языке Lua, что обеспечивает значительную гибкость конфигурирования. Ключевые метрики выходного отчёта~--- количество транзакций в секунду и время отклика. Основным достоинством Sysbench является сочетание простоты установки с широкими возможностями параметризации через Lua-скрипты. Вместе с тем инструмент ограничен интерфейсом командной строки, лишён механизмов хранения истории тестов и не предусматривает автоматизированного сравнения результатов между прогонами, что вынуждает пользователя самостоятельно интерпретировать и сопоставлять текстовые отчёты.

Утилита pgbench входит в стандартную поставку PostgreSQL и моделирует банковскую транзакционную нагрузку по мотивам спецификации TPC-B, измеряя количество транзакций в секунду при заданном числе параллельных клиентских соединений. Инструмент поддерживает выполнение пользовательских SQL-скриптов, что несколько расширяет область его применения за пределы встроенного бенчмарка. Главное достоинство pgbench состоит в отсутствии необходимости установки дополнительного программного обеспечения для пользователей PostgreSQL. Однако жёсткая привязка к одной СУБД, упрощённая модель нагрузки и минимальные возможности визуализации существенно сужают сферу применимости инструмента и делают его скорее отправной точкой для экспресс-оценки, нежели полноценным средством сравнительного анализа производительности.

HammerDB~--- инструмент нагрузочного тестирования и бенчмаркинга баз данных с открытым исходным кодом, разрабатываемый под эгидой организации TPC. Актуальная версия 5.0 распространяется под лицензией GNU GPL v3.0~\cite{hammerdb_v5}. В данном выпуске HammerDB перешёл на бинарные дистрибутивы в виде единых исполняемых файлов для трёх режимов работы: графический интерфейс, интерфейс командной строки и веб-сервис. Ядро инструмента переработано на базе Tcl/Tk\,9.0~--- первого крупного обновления за более чем 27 лет. HammerDB поддерживает MariaDB, PostgreSQL, Oracle, SQL Server, MySQL и IBM\,Db2. Инструмент реализует рабочие нагрузки TPROC-C и TPROC-H~--- производные от спецификаций TPC-C и TPC-H~--- с результирующими метриками NOPM (New Orders Per Minute) и TPM (Transactions Per Minute). Ключевым преимуществом HammerDB является широкая поддержка СУБД в сочетании с соответствием отраслевым стандартам и доступностью официальных образов Docker для интеграции в CI/CD-окружения. Существенным ограничением остаётся отсутствие централизованного механизма хранения истории тестов: каждый прогон представляет собой изолированное событие, а отслеживание динамики изменений производительности возлагается на пользователя. Таким образом, HammerDB закрывает задачу генерации стандартизированной нагрузки, но не решает задачу долгосрочного сравнительного анализа результатов.

\subsubsection{Коммерческие корпоративные решения}

Benchmark Factory 9.1 ~--- коммерческий инструмент тестирования производительности баз данных, позиционируемый как многоплатформенное решение для корпоративного сегмента~\cite{benchmarkfactory_91}. Benchmark Factory позволяет проводить воспроизведение рабочей нагрузки, отраслевые стандартизированные бенчмарки и тестирование масштабируемости. Инструмент доступен для Oracle, SQL Server, IBM DB2, SAP, MySQL, PostgreSQL, MariaDB и других баз данных через ODBC-подключение. Результаты тестирования сохраняются в централизованном репозитории с поддержкой встроенных отчётов и возможностью просмотра нескольких прогонов одновременно; предусмотрен REST API для автоматизации. Ключевое преимущество Benchmark Factory по сравнению с инструментами класса Oracle RAT~--- поддержка нескольких СУБД в рамках единого продукта с GUI. Вместе с тем инструмент является платным решением с графическим интерфейсом настольного приложения, что ограничивает его доступность для академических и небольших проектов и исключает нативный веб-доступ.

Oracle Real Application Testing~--- коммерческий инструмент корпорации Oracle для тестирования производительности баз данных Oracle, доступный исключительно в составе редакции Oracle Enterprise Edition. Компонент Database Replay позволяет зафиксировать реальную рабочую нагрузку продуктивной базы данных и воспроизвести её на тестовой системе, а компонент SQL Performance Analyzer оценивает влияние изменений на выполнение отдельных SQL-запросов, сравнивая планы выполнения и метрики до и после изменений. Несмотря на высокую точность воспроизведения реальной нагрузки, Oracle RAT жёстко привязан к экосистеме Oracle и требует существенных лицензионных затрат, что принципиально ограничивает его применимость в контексте сравнительного тестирования нескольких СУБД.

\subsubsection{Воспроизводимость бенчмаркинга как концептуальный ориентир}

Отдельного упоминания заслуживает инструмент Benchkit~--- открытый фреймворк, решающий не задачу генерации нагрузки, а задачу воспроизводимости бенчмаркинга: декларативное описание окружения, фиксация всех параметров и результатов теста, автоматизация повторных запусков~\cite{benchkit_exasol}. Данная концепция воспроизводимых конвейеров тестирования представляет интерес с точки зрения проектирования системы централизованного хранения и сравнения истории тестов.

Проведённый обзор позволяет констатировать: каждый из рассмотренных инструментов эффективно решает одну конкретную задачу~--- генерацию нагрузки или анализ производительности~--- но ни одно открытое решение не объединяет в едином интерфейсе генерацию нагрузки на реляционные СУБД, мониторинг метрик в реальном времени, централизованное хранение истории тестов и автоматизированное сравнение прогонов. Описанная фрагментированность служит предметом систематического сравнительного анализа в следующем разделе.

\subsection{Сравнительный анализ решений}

Для систематизации результатов обзора проведён сравнительный анализ рассмотренных инструментов по критериям, непосредственно влияющим на удобство и эффективность практики нагрузочного тестирования СУБД. Выбор критериев обусловлен не техническими характеристиками самих инструментов, а практическими требованиями к процессу сравнительного тестирования нескольких СУБД. Пользовательский интерфейс рассматривается как критерий доступности: CLI-инструменты требуют технической экспертизы и не позволяют наблюдать динамику метрик в реальном времени, тогда как браузерный интерфейс снижает порог входа и обеспечивает немедленную визуальную обратную связь без установки дополнительного программного обеспечения. Наличие истории тестов оценивается как критически важная характеристика: разовый прогон фиксирует состояние системы в единственный момент времени, тогда как только накопленная история позволяет выявлять регрессии производительности и долгосрочные тренды после изменений конфигурации или схемы данных. Возможность сравнения результатов производна от наличия истории и определяет, насколько процесс вынесения суждений о производительности автоматизирован, а насколько возложен на пользователя. Результаты анализа сведены в таблицу~\ref{tab:comparison}.

{
\renewcommand{\arraystretch}{0.9}
\setlength{\tabcolsep}{3pt}
\sloppy
\pretolerance=10000
\tolerance=2000
\emergencystretch=1em
\hyphenpenalty=50
\exhyphenpenalty=50

\begin{table}[H]
    \centering
    \caption{Сравнительный анализ инструментов нагрузочного тестирования СУБД}
    \label{tab:comparison}
    \begin{tabularx}{\textwidth}{|>{\RaggedRight}p{0.13\textwidth}|>{\RaggedRight}X|>{\RaggedRight}X|>{\RaggedRight}X|>{\RaggedRight}X|}
        \hline
        Критерий & HammerDB 5.0 & Sysbench / pgbench & Oracle RAT & Benchmark Factory 9.1 \\
        \hline
        Интер\-фейс & GUI\slash CLI\slash Web & CLI & GUI\slash CLI & GUI\slash CLI\slash REST API \\
        \hline
        Под\-дер\-жи\-ва\-емые СУБД & Oracle, SQL Server, Db2, MySQL, MariaDB, PostgreSQL & MySQL, PostgreSQL (pgbench~--- только PostgreSQL) & Только Oracle Database & Oracle, SQL Server, Db2, MySQL, PostgreSQL, SAP, MariaDB и др. \\
        \hline
        Хранение истории тестов & Отсутствует (логи) & Отсутствует & Частично & Цент\-ра\-ли\-зо\-ван\-ный репо\-зи\-то\-рий \\
        \hline
        Сравнение тестов & Ручной анализ логов & Ручное & Встроенное & Встроенные отчёты \\
        \hline
        Ви\-зу\-а\-ли\-за\-ция & Простые графики & Текстовый отчёт & Детальные отчёты OEM & Встроенные отчёты \\
        \hline
        Модель рас\-про\-стра\-не\-ния & Открытая & Открытая & Прo\-прие\-тар\-ная & Прo\-прие\-тар\-ная \\
        \hline
    \end{tabularx}
\end{table}
}

Анализ таблицы~\ref{tab:comparison} выявляет устойчивые закономерности, характерные для каждого класса инструментов, которые не сводятся к отдельным достоинствам и недостаткам конкретных продуктов.

Инструменты с открытым исходным кодом~--- HammerDB, Sysbench и pgbench~--- демонстрируют общий паттерн: высокая гибкость в части генерации нагрузки при принципиальном отсутствии механизмов долгосрочного накопления и сравнения результатов. Каждый прогон существует как изолированный артефакт; сопоставление результатов между прогонами остаётся ручной операцией, полностью возлагаемой на пользователя. Такая архитектура оправдана для разовых замеров производительности, однако делает перечисленные инструменты непригодными для систематического отслеживания изменений производительности во времени без создания вспомогательной инфраструктуры хранения и анализа данных.

Коммерческие продукты~--- Oracle RAT и Benchmark Factory 9.1~--- предлагают наиболее развитый функционал в части анализа: централизованные репозитории результатов, встроенные сравнительные отчёты, поддержку нескольких типов рабочей нагрузки. Однако данный класс инструментов ограничен по-разному: Oracle RAT жёстко привязан к экосистеме одного поставщика и недоступен вне Oracle Enterprise Edition, тогда как Benchmark Factory, при широте поддерживаемых СУБД, является платным продуктом с настольным графическим интерфейсом, что исключает нативный веб-доступ и существенно ограничивает применимость в академических и некоммерческих проектах.

Таким образом, рассмотренные инструменты образуют фрагментированную экосистему, в которой ни одно решение с открытым исходным кодом не обеспечивает одновременно генерацию нагрузки на несколько реляционных СУБД, мониторинг метрик в реальном времени через браузерный интерфейс, централизованное хранение истории тестов и автоматизированное сравнение прогонов. Коммерческие продукты частично закрывают этот разрыв, однако за счёт закрытости экосистемы или высокой стоимости лицензирования. Выявленный функциональный разрыв определяет состав требований к разрабатываемой системе, сформулированных в следующем разделе.

\subsection{Требования к разрабатываемой системе}

На основании проведённого аналитического обзора и систематического выявления структурных недостатков существующих решений сформулированы требования к разрабатываемой системе сравнительного нагрузочного тестирования реляционных СУБД. Требования разделены на функциональные, нефункциональные и ограничения; каждое из них непосредственно вытекает из функционального разрыва, зафиксированного в разделе~1.3.

Функциональные требования определяют действия, которые система обязана выполнять. Система должна обеспечивать поддержку трёх реляционных СУБД~--- MySQL, MariaDB и PostgreSQL~--- с возможностью сохранения нескольких конфигураций подключений и их совместного использования в рамках одного тестового прогона. Данное требование вытекает из необходимости обеспечить единую точку управления разнородными СУБД, что является ключевым пробелом существующих открытых инструментов.

Система должна поддерживать объединение нескольких физических подключений с одинаковой структурой данных в группу и автоматически подбирать к ней подходящий набор сценариев нагрузки, что позволяет проводить сравнительное тестирование без ручного конфигурирования каждого подключения. Генерация нагрузки должна осуществляться по параметризованным сценариям, которые задают состав операций~--- как отдельных SQL-запросов, так и последовательностей операций в рамках атомарных транзакций~--- с настраиваемым соотношением типов, числом виртуальных пользователей, числом итераций и длительностью прогрева.

Система должна обеспечивать потоковую передачу метрик производительности в веб-интерфейс с задержкой не более одной секунды и параллельно сохранять историю прогонов с возможностью фильтрации и просмотра ранее выполненных тестов. Система должна предоставлять инструменты сравнительного анализа двух и более прогонов~--- как попарного сравнения СУБД при одинаковой нагрузке, так и анализа изменения характеристик при последовательном увеличении числа виртуальных пользователей~--- с автоматическим применением статистических критериев, вычислением размера эффекта и формированием структурированного отчёта.

Система должна автоматически сохранять состояние тестируемых баз данных перед прогонами, включающими операции записи, и восстанавливать его по завершении теста с верификацией целостности данных, что обеспечивает воспроизводимость экспериментов.

Нефункциональные требования определяют качественные характеристики системы. Веб-интерфейс должен быть доступен через стандартный браузер без установки дополнительного программного обеспечения на стороне клиента~--- в противовес настольным GUI-решениям, характерным для коммерческого сегмента. Архитектура системы должна допускать расширение перечня поддерживаемых СУБД без значительной модификации кодовой базы, обеспечивая открытость к развитию. Развёртывание должно осуществляться с использованием контейнеризации, что гарантирует воспроизводимость тестового окружения и упрощает установку в различных инфраструктурах.

Ограничения фиксируют рамки применимости системы в текущей версии. Система предназначена для сравнительного анализа производительности, а не для проведения формальных сертифицированных бенчмарков TPC. Генерация нагрузки осуществляется на одном сервере; горизонтальное масштабирование генераторов нагрузки в текущей версии не предусматривается. Область применения системы ограничена реляционными СУБД.

Сформулированные требования определяют функциональные и технические характеристики разрабатываемой системы и составляют непосредственную основу для проектирования её архитектуры, рассматриваемого во второй главе.

\subsection{Выводы по главе}

Анализ существующих инструментов нагрузочного тестирования показал, что ни одно открытое решение не сочетает генерацию нагрузки на несколько СУБД, мониторинг метрик в реальном времени через браузер, централизованное хранение истории тестов и автоматизированный сравнительный анализ со статистикой. Выявленный пробел позволил сформулировать восемь функциональных требований (включая поддержку трёх СУБД в едином интерфейсе, группы подключений, параметризованные сценарии, потоковую передачу метрик, историю прогонов, двухрежимный анализ и автоматическое резервное копирование), а также нефункциональные требования и ограничения.